% P4-b — Nueva estrategia de descomposición. Sin preámbulo ni \subsection.

Richards y Ford (2020, cap.~17) proponen tres guías para encontrar la frontera
correcta de un servicio, y las tres fallaron en el diseño original porque nunca se
aplicaron: solo se usó la estructura de la base de datos.

\paragraph{Propósito.}
Cada servicio debe ser funcionalmente cohesivo en torno a un comportamiento
significativo del dominio, no un contenedor de operaciones CRUD sobre una tabla
(Richards \& Ford, 2020, cap.~17). Un servicio ``Empleado'' que solo expone altas,
bajas y modificaciones del registro no tiene un propósito de negocio propio: es la
base de datos con una interfaz de red delante.

\paragraph{Transacciones.}
Las entidades que deben cambiar juntas dentro de la misma transacción de negocio
a menudo señalan una buena frontera de servicio (Richards \& Ford, 2020,
cap.~17). Suponemos que contratar a un empleado es una única transacción de
negocio --- que el alta del empleado, la configuración inicial de su nómina y
su inscripción en el plan de beneficios deben confirmarse o revertirse
juntas ---, una decisión que debe validarse con el dueño del dominio y no una
propiedad universal del proceso. Bajo ese supuesto, esas tres partes deben
vivir dentro del mismo servicio, en vez de repartirse entre tres servicios que
luego necesitan coordinarse por la red.

\paragraph{Coreografía.}
Cuando un conjunto de servicios con buen aislamiento de dominio termina
necesitando comunicación excesiva entre sí para completar una operación, el
arquitecto puede considerar reagruparlos en un servicio más grande, en vez de
multiplicar las llamadas entre ellos (Richards \& Ford, 2020, cap.~17). Esa es exactamente la situación de
RRHH: el aislamiento por entidad parecía limpio en el papel, pero el costo real
apareció al ejecutar el \emph{workflow} completo de contratación.

Aplicando las tres guías, la nueva frontera deja de ser la tabla de la base de
datos y pasa a ser el \emph{workflow} de contratación: Empleado, la configuración
inicial de Nómina y la inscripción en Beneficios se agrupan en un único
\texttt{ServicioContratación}, porque comparten el mismo propósito (dar de alta a
un empleado) y la misma frontera transaccional (deben confirmarse o revertirse
como una unidad). Departamento, en cambio, no participa en esa transacción ---
solo se consulta como dato de referencia para validar a qué equipo se asigna el
nuevo empleado ---, así que puede seguir siendo un servicio independiente sin
introducir una transacción distribuida: la consulta no necesita ser atómica con
el alta.

Todo lo que no forma parte de ese núcleo transaccional se resuelve por
\textbf{coreografía asíncrona basada en eventos}, en vez de llamadas síncronas
encadenadas: aprovisionar el equipo de TI, notificar al gerente del área, actualizar el directorio corporativo. En vez de
que \texttt{ServicioContratación} llame uno por uno a esos servicios y espere su
respuesta, publica un único evento \texttt{EmpleadoContratado} al confirmar la
transacción, y cada servicio periférico reacciona de forma asíncrona y a su propio
ritmo. Esto reduce la coreografía síncrona obligatoria de ``casi todos los
servicios'' a una sola llamada atómica dentro de una frontera transaccional bien
definida, más un evento que dispara el resto del proceso sin bloquear la
confirmación de la contratación.

\begin{table}[H]
  \centering
  \caption{Nueva estrategia de descomposición para la operación de contratación.}
  \label{tab:p4b-comparacion}
  \small
  \renewcommand{\arraystretch}{1.35}
  \begin{tabularx}{\textwidth}{>{\raggedright\arraybackslash}p{2.7cm} >{\raggedright\arraybackslash}X >{\raggedright\arraybackslash}X}
    \toprule
    \textbf{Criterio} & \textbf{Antes (por entidad)} & \textbf{Después (por workflow)} \\
    \midrule
    Propósito &
      un servicio por tabla (Empleado, Departamento, Nómina, Beneficios) &
      un servicio por transacción de negocio (\texttt{ServicioContratación}) \\
    Transacciones &
      ``contratar'' cruza síncronamente los 4 servicios &
      ``contratar'' es una sola llamada atómica dentro de \texttt{ServicioContratación} \\
    Dato de referencia &
      Departamento consultado síncronamente en cada paso &
      Departamento se consulta fuera de la transacción, como referencia \\
    Efectos secundarios &
      notificar a TI, gerente y directorio: llamadas síncronas adicionales &
      evento \texttt{EmpleadoContratado}; TI, gerente y directorio reaccionan por
      coreografía \\
    \bottomrule
  \end{tabularx}
\end{table}

El costo de esta estrategia es el mismo que ya se discutió para Service-Based
Architecture en P1: \texttt{ServicioContratación} concentra más responsabilidad y
más riesgo de despliegue que los cuatro microservicios originales. Pero es un
costo deliberado: se paga acoplamiento interno, local y controlado, a cambio de
eliminar la transacción distribuida que hacía frágil al diseño anterior --- el
mismo trade-off de siempre, ahora aplicado a la granularidad en vez de a la
comunicación.
